% Chapter V, Section 50 (first half, equations 1--10)
% Chandrasekhar, "Hydrodynamic and Hydromagnetic Stability"
% Book pages 197--198

\subsection*{\S\,50. The propagation of hydromagnetic waves in a rotating fluid}

In Chapters~III and~IV we saw that the propagation of waves in a
rotating fluid and in the presence of a magnetic field bear on the respective
problems in thermal instability.  We shall see that the propagation
of hydromagnetic waves in a rotating fluid has a similar bearing on the
problem to be investigated in this chapter.

Consider an incompressible fluid in a state of uniform rotation with
an angular velocity~$\boldsymbol{\Omega}$.  In a frame of reference rotating with the angular
velocity~$\boldsymbol{\Omega}$, the basic equations are (cf.\ equations III~(44) and IV~(10)):
\begin{align}
  \frac{\partial u_i}{\partial t}
  + u_j\,\frac{\partial u_i}{\partial x_j}
  - \frac{\mu H_j}{4\pi\rho}\,\frac{\partial H_i}{\partial x_j}
  \notag\\
  &= 2\epsilon_{ijk}\,u_j\,\Omega_k
    - \frac{\partial}{\partial x_i}\!\left(
        p + \frac{|\mathbf{H}|^{2}}{8\pi}
        - \tfrac{1}{2}|\boldsymbol{\Omega}\times\mathbf{r}|^{2} + V
    \right),
  \label{eq:5-1}
\end{align}
\begin{equation}
  \frac{\partial H_i}{\partial t}
  + u_j\,\frac{\partial H_i}{\partial x_j}
  - H_j\,\frac{\partial u_i}{\partial x_j}
  = \eta\,\nabla^{2}H_i,
  \label{eq:5-2}
\end{equation}
\begin{equation}
  \frac{\partial u_i}{\partial x_i} = 0,
  \quad\text{and}\quad
  \frac{\partial H_i}{\partial x_i} = 0.
  \label{eq:5-3}
\end{equation}

For an inviscid fluid of zero resistivity, the equations governing small
departures from an initial state in which a uniform magnetic field~$H_j$
prevails are
\begin{equation}
  \frac{\partial u_i}{\partial t}
  - \frac{\mu H_j}{4\pi\rho}\,\frac{\partial h_i}{\partial x_j}
  = -\frac{\partial}{\partial x_i}\,\delta\varpi
    + 2\epsilon_{ijk}\,u_j\,\Omega_k,
  \label{eq:5-4}
\end{equation}
\begin{equation}
  \frac{\partial h_i}{\partial t}
  = H_j\,\frac{\partial u_i}{\partial x_j},
  \quad
  \frac{\partial u_i}{\partial x_i} = 0,
  \quad\text{and}\quad
  \frac{\partial h_i}{\partial x_i} = 0,
  \label{eq:5-5}
\end{equation}
where~$h_i$ denotes the departure of the magnetic field from~$H_i$.

We now seek solutions of equations~\eqref{eq:5-4} and~\eqref{eq:5-5}
whose space-time dependence is given by
\begin{equation}
  e^{i(\sigma t + \mathbf{k}\cdot\mathbf{x})}.
  \label{eq:5-6}
\end{equation}

Equations~\eqref{eq:5-4} and~\eqref{eq:5-5} then give
\begin{equation}
  \sigma u_i
  - \frac{\mu(H_j k_j)}{4\pi\rho}\,h_i
  = -k_i\,\delta\varpi - 2i\epsilon_{ijk}\,u_j\,\Omega_k,
  \label{eq:5-7}
\end{equation}
\begin{equation}
  \sigma h_i = (H_j k_j)\,u_i,
  \label{eq:5-8}
\end{equation}
\begin{equation}
  u_i\,k_i = 0,
  \quad\text{and}\quad
  h_i\,k_i = 0.
  \label{eq:5-9}
\end{equation}

Equation~\eqref{eq:5-9} shows that \emph{the waves are transverse}.

Eliminating~$h_i$ between equations~\eqref{eq:5-7} and~\eqref{eq:5-8},
we obtain
\begin{equation}
  n u_i + k_i\,\delta\varpi + 2i\epsilon_{ijk}\,u_j\,\Omega_k = 0,
  \label{eq:5-10}
\end{equation}
